% Copyright 2026 Yoshida Shin
%
% This file is part of the ``Repeating Nines''.
%
% Permission is granted to copy, distribute, and/or modify this document
% under the terms of the GNU Free Documentation License, Version 1.3 or any later version published by the Free Software Foundation;
% with no Front-Cover Texts, no Back-Cover Texts, and one Invariant Section: ``This document was written by Shin Yoshida with the aim of contributing to a fairer and safer world.''
% A copy of the license is included in the section entitled ``GNU Free Documentation License.''


\begin{frame}[c]{}{}
  {\large Rigorous Definition}
\end{frame}

\begin{frame}[c]{}{}
  Imagine a frivolous game using $a_n$ between two players, A and B.
\end{frame}

\begin{frame}[c]{}{}
  First Move:

  \vspace{4ex}

  A: State a real number $k_1$

  \vspace{2ex}

  B: If B can find such an n

  $$n \in \mathbb{N}, k_1 \leqq a_n$$

  the game ends and B wins!
\end{frame}

\begin{frame}[c]{}{}
  Second Move:

  \vspace{4ex}

  B: State a real number $k_2 (< k_1)$

  \vspace{2ex}

  A: If A can find such an n

  $$n \in \mathbb{N}, k_2 \leqq a_n$$

  the game ends and A wins!
\end{frame}

\begin{frame}[c]{}{}
  Third Move:

  \vspace{4ex}

  A: State a real number $k_3 (< k_2)$

  ...
\end{frame}

\begin{frame}[c]{}{}
  Does the rule make sense?

  \vspace{4ex}

  For example, the game might go as follows:
\end{frame}

\begin{frame}[c]{}{}
  First Move:

  \vspace{4ex}

  A: $k_1 = \pi$

  ($\pi$ is a real number, it is OK.)

  \vspace{2ex}

  B: (I can't find such an n $\pi \leqq a_n$)
\end{frame}

\begin{frame}[c]{}{}
  Second Move:

  \vspace{4ex}

  B: $k_2$ = 1.5

  $(1.5$ is a real number, that is less than $k_1 = \pi$.)

  \vspace{2ex}

  A: (I can't find such an n $1.5 \leqq a_n$)
\end{frame}

\begin{frame}[c]{}{}
  Third Move:

  \vspace{4ex}

  A: $k_3$ = -5

  (-5 is a real number, that is less than $k_2$ = 1.5.)

  \vspace{2ex}

  B: $n = 3$

  ($a_3 = 0.999, -5 \leqq a_3$, I won!)
\end{frame}

\begin{frame}[c]{}{}
  Do you understand the rules?
\end{frame}

\begin{frame}[c]{}{}
  Then, let's consider why is this game ``frivolous.''
\end{frame}

\begin{frame}[c]{}{}
  Because player A can always win if they choose the right number in their first move.

  \vspace{4ex}

  Let's call that number ``0.9999...''
\end{frame}

\begin{frame}[c]{}{}
  In other words, ``0.9999...'' is a value fulfilling the following conditions.

  \vspace{2ex}

  \begin{itemize}
  \item $\forall n \in \mathbb{N}, a_n < 0.9999...$
  \item $\forall k < 0.9999..., \exists n$ $k \leqq a_n$
  \end{itemize}

  \vspace{2ex}

  If and only if A states such 0.9999... at first,

  A can always win against B.
\end{frame}

\begin{frame}[c]{}{}
  Note that, I have not proved that

  ``One, or more than one value fulfills those conditions.''

  \vspace{2ex}

  i.e. No such a value may exist!
\end{frame}
